% 来源：OEBPS/Text/01_28.xhtml（OPF spine 顺序；仅第一册）
\cleardoublepage
\BookEntry[本杰明·M.弗里德曼《经济增长的道德意义》]{本杰明·M.弗里德曼《经济增长的道德意义》}

\BookCover{images/书单14经济增长的道德意义.jpg}{[美] 本杰明·M.弗里德曼著| 李天有译| 中国人民大学出版社2013年版}

\begin{BookSummaryBox}
\textbf{具有跨学科思想的经济理论著作}
\par
\end{BookSummaryBox}

\subsubsection*{这本书谈什么？}

对任何一个国家，经济增长都是一个古老的话题。在《经济增长的道德意义》中，弗里德曼开门见山地提出：“我们那么明确而极大地关注经济增长是否正确呢？”在做了大量细致的跨学科研究之后，弗里德曼最终指出：增长，而不仅仅是生活水平，是在第三世界政治与社会自由化的关键。但是，他也警告说，即便是像美国这样富裕的国家，一旦收入增长长期停滞，民主价值也会处于危险之中。
\par

\subsubsection*{芒格为什么谈起这本书？}

芒格在发表了“论学院派经济学”的演讲之后，看到一本书，是阿尔弗雷德·诺普夫出版社在2005年出版的。芒格说：
\par

\begin{quotation}
它的作者是杰出的哈佛大学经济学教授本杰明·弗里德曼。如我在演讲中所希望的那样，这本书讨论了经济学和道德之间的相互关系。这本书的名字是《经济增长的道德意义》。读者从书名可以看出，弗里德曼教授尤其感兴趣的是经济增长对道德的影响，而我感兴趣的则与此相反，主要是道德对经济增长的影响。这个区别关系不大，因为任何受过教育的人都能明白这两者会相互产生或好或坏的影响，造成通常所说的“良性循环”或者“恶性循环”。弗里德曼教授给这个主题增添了一句来自艾利沙·本·阿萨里亚拉比（公元前1世纪左右第二代犹太密西拿学者）的名言：“没有面包就没有法律，没有法律就没有面包。”
\par
\end{quotation}

\subsubsection*{《经济增长的道德意义》作者本杰明·M.弗里德曼如是说}

“关于经济增长的问题并不是一个物质或道德价值的问题。确实，经济增长常常伴随着不必要的副作用，如对传统文化的破坏和对环境的损害，但是这些问题中有一些确实是我们有权利加以考察的一种正当的道德关怀。虽然如此，经济增长也包含着道德上有益的社会和政治结果。特别是在评价不同的公共政策时，重要的是不仅要考虑我们熟悉的道德上的负面因素，同时也要考虑到这些道德上的正面因素。”
\par

“正确理解有关经济增长的公共政策和私人行动之间的关系也一样至关重要。这里，同样的结果是，上升的生活水准的正面道德结果会显著地改编故事的内容。”
\par
